\documentclass[10pt,a4paper]{amsart}
\usepackage{amsmath,amssymb,amsthm,graphicx}
% Retained minimal preamble for the Stacks Project chapter more-algebra.tex (More on Algebra).
% The chapter file pinned-original-source.tex carries only \input{preamble} on its first line, so no
% preamble text lives inside the chapter itself.  The Stacks Project uses one single shared preamble for
% every chapter; the snapshot stages it identically as sources/stacks-chow/preamble.tex and
% sources/stacks-etale-cohomology/preamble.tex.  The block below is copied mechanically, line for line,
% from sources/stacks-chow/preamble.tex lines 14-18 (the Xy-pic setup) and 154-217 (the native theorem
% declarations together with the mathematical macros).  Nothing was added, removed or redefined; the
% wrapper class and the amsmath/amssymb/amsthm/graphicx base packages are supplied by the assembler.
% Chapter: more-algebra.tex -- More on Algebra.
% For commutative diagrams we use Xy-pic
\usepackage[all]{xy}

% We use 2cell for 2-commutative diagrams.
\xyoption{2cell}

% Theorem environments.
%
\theoremstyle{plain}
\newtheorem{theorem}[subsection]{Theorem}
\newtheorem{proposition}[subsection]{Proposition}
\newtheorem{lemma}[subsection]{Lemma}

\theoremstyle{definition}
\newtheorem{definition}[subsection]{Definition}
\newtheorem{example}[subsection]{Example}
\newtheorem{exercise}[subsection]{Exercise}
\newtheorem{situation}[subsection]{Situation}

\theoremstyle{remark}
\newtheorem{remark}[subsection]{Remark}
\newtheorem{remarks}[subsection]{Remarks}

\numberwithin{equation}{subsection}

% Macros
%
\def\lim{\mathop{\mathrm{lim}}\nolimits}
\def\colim{\mathop{\mathrm{colim}}\nolimits}
\def\Spec{\mathop{\mathrm{Spec}}}
\def\Hom{\mathop{\mathrm{Hom}}\nolimits}
\def\Ext{\mathop{\mathrm{Ext}}\nolimits}
\def\SheafHom{\mathop{\mathcal{H}\!\mathit{om}}\nolimits}
\def\SheafExt{\mathop{\mathcal{E}\!\mathit{xt}}\nolimits}
\def\Sch{\mathit{Sch}}
\def\Mor{\mathop{\mathrm{Mor}}\nolimits}
\def\Ob{\mathop{\mathrm{Ob}}\nolimits}
\def\Sh{\mathop{\mathit{Sh}}\nolimits}
\def\NL{\mathop{N\!L}\nolimits}
\def\CH{\mathop{\mathrm{CH}}\nolimits}
\def\proetale{{pro\text{-}\acute{e}tale}}
\def\etale{{\acute{e}tale}}
\def\QCoh{\mathit{QCoh}}
\def\Ker{\mathop{\mathrm{Ker}}}
\def\Im{\mathop{\mathrm{Im}}}
\def\Coker{\mathop{\mathrm{Coker}}}
\def\Coim{\mathop{\mathrm{Coim}}}

% Boxtimes
%
\DeclareMathSymbol{\boxtimes}{\mathbin}{AMSa}{"02}

%
% Macros for moduli stacks/spaces
%
\def\QCohstack{\mathcal{QC}\!\mathit{oh}}
\def\Cohstack{\mathcal{C}\!\mathit{oh}}
\def\Spacesstack{\mathcal{S}\!\mathit{paces}}
\def\Quotfunctor{\mathrm{Quot}}
\def\Hilbfunctor{\mathrm{Hilb}}
\def\Curvesstack{\mathcal{C}\!\mathit{urves}}
\def\Polarizedstack{\mathcal{P}\!\mathit{olarized}}
\def\Complexesstack{\mathcal{C}\!\mathit{omplexes}}
% \Pic is the operator that assigns to X its picard group, usage \Pic(X)
% \Picardstack_{X/B} denotes the Picard stack of X over B
% \Picardfunctor_{X/B} denotes the Picard functor of X over B
\def\Pic{\mathop{\mathrm{Pic}}\nolimits}
\def\Picardstack{\mathcal{P}\!\mathit{ic}}
\def\Picardfunctor{\mathrm{Pic}}
\def\Deformationcategory{\mathcal{D}\!\mathit{ef}}

% --- appended 2026-10-12 (Chief): the project own definitions of the reference/slogan/history
% display environments, copied mechanically from the shared Stacks preamble lines 8-12
% (\usepackage{verbatim} plus the three \newenvironment lines, which the chapter preamble keeps
% verbatim from the project preamble). Nothing added, removed or redefined: these environments
% carry no mathematical content and the project itself typesets them as comments. Without these
% three lines a counted statement whose own text uses them fails with Environment slogan undefined.
\usepackage{verbatim}
\newenvironment{reference}{\comment}{\endcomment}
\newenvironment{slogan}{\comment}{\endcomment}
\newenvironment{history}{\comment}{\endcomment}
\usepackage{cleveref}
\crefname{originallemma}{original lemma}{original lemmas}
\Crefname{originallemma}{Original lemma}{Original lemmas}
\makeatletter
\newlabel{more-algebra-proposition-finite-type-over-P-ring}{{0BIV}{source}{}{sourcepointer-1}{}}
\newlabel{more-algebra-proposition-finite-type-over-P-ring@cref}{{[originallemma][0][]0BIV}{[1][0][]source}{}{}{}}
\makeatother
\crefname{originallemma}{original lemma}{original lemmas}
\Crefname{originallemma}{Original lemma}{Original lemmas}
\makeatletter
\newlabel{lemma-P-ring-goes-up-quasi-finite}{{0BIT}{source}{}{sourcepointer-2}{}}
\newlabel{lemma-P-ring-goes-up-quasi-finite@cref}{{[originallemma][0][]0BIT}{[1][0][]source}{}{}{}}
\makeatother
\crefname{originallemma}{original lemma}{original lemmas}
\Crefname{originallemma}{Original lemma}{Original lemmas}
\makeatletter
\newlabel{lemma-check-P-ring-maximal-ideals}{{0BIU}{source}{}{sourcepointer-3}{}}
\newlabel{lemma-check-P-ring-maximal-ideals@cref}{{[originallemma][0][]0BIU}{[1][0][]source}{}{}{}}
\makeatother
\crefname{originallemma}{original lemma}{original lemmas}
\Crefname{originallemma}{Original lemma}{Original lemmas}
\makeatletter
\newlabel{lemma-P-local}{{0BK8}{source}{}{sourcepointer-4}{}}
\newlabel{lemma-P-local@cref}{{[originallemma][0][]0BK8}{[1][0][]source}{}{}{}}
\makeatother
\crefname{originallemma}{original lemma}{original lemmas}
\Crefname{originallemma}{Original lemma}{Original lemmas}
\makeatletter
\newlabel{proposition-Noetherian-complete-G-ring}{{07PS}{source}{}{sourcepointer-5}{}}
\newlabel{proposition-Noetherian-complete-G-ring@cref}{{[originallemma][0][]07PS}{[1][0][]source}{}{}{}}
\makeatother
\crefname{originallemma}{original lemma}{original lemmas}
\Crefname{originallemma}{Original lemma}{Original lemmas}
\makeatletter
\newlabel{proposition-finite-type-over-G-ring}{{07PV}{source}{}{sourcepointer-6}{}}
\newlabel{proposition-finite-type-over-G-ring@cref}{{[originallemma][0][]07PV}{[1][0][]source}{}{}{}}
\makeatother
\title{Essentially finite type algebras over a P-ring are P-rings}
\author{The Stacks Project Authors}
\date{stacks-project:more-algebra.tex:local-snapshot-20261012; GNU FDL 1.2 or later}
\begin{document}
\maketitle
\setcounter{subsection}{1}
\noindent\textit{Editorial note.} All delivered mathematical text is the author's own, copied from the pinned Stacks Project chapter file more-algebra.tex; no statement or proof step was rewritten, reordered or removed.\par\medskip
\noindent\textit{Editorial note.} The chapter file pinned-original-source.tex carries only \textbackslash{}input\{preamble\} on its first line, so no preamble text lives inside the chapter. The retained file sources/stacks-more-algebra/retained-preamble.tex is a mechanical, line-for-line copy of the project's single shared preamble as staged in the snapshot (Xy-pic setup plus theorem declarations and mathematical macros). Nothing was added, removed or redefined, and the excerpt serves only as preamble, never as a span coordinate.\par\medskip
\noindent\textit{Editorial note.} Both counted spans start at the \textbackslash{}begin line of their environment, so the statement and the author's proof are delivered whole; surrounding prose of the section is deliberately outside the counted spans and is carried as standard prerequisites instead.\par\medskip
\noindent\textit{Editorial note.} The citation fields (source\_id, source\_version, authors, source\_url, license.name, source\_title) are identical to those of every other card this lane has produced from this chapter, so the batch selector treats them as one source.\par\medskip
\noindent\textit{Editorial note.} Benign-screen note: the proof contains the phrase 'it is clear that' / 'is obvious' once, used as a legitimate step inside a complete argument; the author proves the statement, so this is not a hint-only proof.\par\medskip
\noindent\textit{Allowed standard prerequisites (editorial).}\begin{itemize}
\item Being a P-ring is local on the ring, and it suffices to check the formal fibres at maximal ideals. Lemma \textbackslash{}ref\{lemma-check-P-ring-maximal-ideals\} and Lemma \textbackslash{}ref\{lemma-P-local\}, cited at proof lines 13015 and 13032-13033. Reduce the claim to a statement about formal fibres.
\item Ascent of the P-ring property along quasi-finite maps, and the P-ring property of polynomial rings. Lemma \textbackslash{}ref\{lemma-P-ring-goes-up-quasi-finite\}, cited at proof lines 13010-13011. Reduce to the polynomial algebra case.
\item Complete Noetherian rings are G-rings, and finite type algebras over G-rings are G-rings. Propositions \textbackslash{}ref\{proposition-Noetherian-complete-G-ring\} and \textbackslash{}ref\{proposition-finite-type-over-G-ring\}, cited at proof lines 13036-13038. Make the right vertical arrow of the comparison diagram regular.
\item The permanence conditions (A), (C) and (D) for property P. used at proof lines 13029-13031, 13039-13041, and 13042-13044. Carry the fibrewise property along the two arrows of the diagram.
\end{itemize}\medskip
\ifdefined\hypertarget\hypertarget{sourcepointer-1}{}\fi
\noindent\textit{Source reference.} Original lemma 0BIV: focused-seeds/AGP/sources/stacks-snapshot-20261011/tags.txt, source lines 14503--14503. The lemma's own canonical tag, quoted so that readers can resolve the label inside the delivered text; it points at the counted statement itself and carries no proof content.\par\medskip
\ifdefined\hypertarget\hypertarget{sourcepointer-2}{}\fi
\noindent\textit{Source reference.} Original lemma 0BIT: focused-seeds/AGP/sources/stacks-snapshot-20261011/tags.txt, source lines 14501--14501. Navigation pointer to a statement the counted proof refers to but does not contain; nothing load-bearing is delegated to it.\par\medskip
\ifdefined\hypertarget\hypertarget{sourcepointer-3}{}\fi
\noindent\textit{Source reference.} Original lemma 0BIU: focused-seeds/AGP/sources/stacks-snapshot-20261011/tags.txt, source lines 14502--14502. Navigation pointer to a statement the counted proof refers to but does not contain; nothing load-bearing is delegated to it.\par\medskip
\ifdefined\hypertarget\hypertarget{sourcepointer-4}{}\fi
\noindent\textit{Source reference.} Original lemma 0BK8: focused-seeds/AGP/sources/stacks-snapshot-20261011/tags.txt, source lines 14551--14551. Navigation pointer to a statement the counted proof refers to but does not contain; nothing load-bearing is delegated to it.\par\medskip
\ifdefined\hypertarget\hypertarget{sourcepointer-5}{}\fi
\noindent\textit{Source reference.} Original lemma 07PS: focused-seeds/AGP/sources/stacks-snapshot-20261011/tags.txt, source lines 9807--9807. Navigation pointer to a statement the counted proof refers to but does not contain; nothing load-bearing is delegated to it.\par\medskip
\ifdefined\hypertarget\hypertarget{sourcepointer-6}{}\fi
\noindent\textit{Source reference.} Original lemma 07PV: focused-seeds/AGP/sources/stacks-snapshot-20261011/tags.txt, source lines 9810--9810. Navigation pointer to a statement the counted proof refers to but does not contain; nothing load-bearing is delegated to it.\par\medskip
\begin{proposition}
\label{proposition-finite-type-over-P-ring}
Let $R$ be a $P$-ring where $P$ satisfies (A), (B), (C), and (D).
If $R \to S$ is essentially of finite type then $S$ is a $P$-ring.
\end{proposition}

\begin{proof}
Since being a $P$-ring is a property of the local rings it is clear
that a localization of a $P$-ring is a $P$-ring. Conversely, if every
localization at a prime is a $P$-ring, then the ring is a $P$-ring.
Thus it suffices to show that $S_\mathfrak q$ is a $P$-ring for every
finite type $R$-algebra $S$ and every prime $\mathfrak q$ of $S$.
Writing $S$ as a quotient of $R[x_1, \ldots, x_n]$ we see from
Lemma \ref{lemma-P-ring-goes-up-quasi-finite} that it suffices to prove
that $R[x_1, \ldots, x_n]$ is a $P$-ring. By induction on $n$ it
suffices to prove that $R[x]$ is a $P$-ring. Let $\mathfrak q \subset R[x]$
be a maximal ideal. By Lemma \ref{lemma-check-P-ring-maximal-ideals}
it suffices to show that the fibres of
$$
R[x]_\mathfrak q \longrightarrow R[x]_\mathfrak q^\wedge
$$
have $P$. If $\mathfrak q$ lies over $\mathfrak p \subset R$, then
we may replace $R$ by $R_\mathfrak p$. Hence we may assume that $R$
is a Noetherian local $P$-ring with maximal ideal $\mathfrak m$ and
that $\mathfrak q \subset R[x]$ lies over $\mathfrak m$. Note that
there is a unique prime $\mathfrak q' \subset R^\wedge[x]$
lying over $\mathfrak q$. Consider the diagram
$$
\xymatrix{
R[x]_\mathfrak q^\wedge \ar[r] &
(R^\wedge[x]_{\mathfrak q'})^\wedge \\
R[x]_\mathfrak q \ar[r] \ar[u] & R^\wedge[x]_{\mathfrak q'} \ar[u]
}
$$
Since $R$ is a $P$-ring the fibres of $R[x] \to R^\wedge[x]$ have
$P$ because they are base changes of the fibres of $R \to R^\wedge$
by a finitely generated field extension so (A) applies. Hence
the fibres of the lower horizontal arrow have $P$ for example by
Lemma \ref{lemma-P-local}.
The right vertical arrow is regular because $R^\wedge$ is
a G-ring (Propositions \ref{proposition-Noetherian-complete-G-ring} and
\ref{proposition-finite-type-over-G-ring}).
It follows that the fibres of the composition
$R[x]_\mathfrak q \to (R^\wedge[x]_{\mathfrak q'})^\wedge$
have $P$ by (C). Hence
the fibres of the left vertical arrow have $P$ by (D) and the
proof is complete.
\end{proof}

\end{document}
